%%%%%
%
\section{D�riv�e au 1er ordre d'un angle, en radian ou en degr�s}

\itemstyle{Op\'erateur~:} \operfmt{Deriv_Angl_Univ}
\index{Deriv_Angl_Univ}

%%%%%OBJET%%%%
\itemstyle{Objet~:}
Calcule la d�riv�e centr�e discr�te, par rapport au temps, d'une variable 
angulaire de type \texttt{float} �chantillonn�e r�guli�rement,
et exprim�e en degr�s ou en radians.

%%%%%ARGUMENTS%%%%
\bigskip
%%%%%ENTREES%%%%
\begin{tabarg}
\multicolumn{3}{l}{\titretab {Entr\'ees}}\\
\hline
\$Arg1 & Vecteur, {\tt float} & blabla (any)\\
\$Arg2 & Coeff., {\tt float} & blabla (any)\\
\$Arg3 & Coeff., {\tt String         } & blabla ()\\
\end{tabarg}

\medskip
%%%%%SORTIES%%%%
\begin{tabarg}
\multicolumn{3}{l}{\titretab {Sorties}}\\
\hline
\$Arg4 & Vecteur, {\tt float} & blabla ($Arg1.unit+ ".s-1" )\\
\end{tabarg}


%%%%%FORMULE%%%%
\itemstyle{Formule}

%%Deriv_Angl_Univ
\[
Resultat = a-b. \Zeta
\]



%%%%%AUTEUR%%%%
\itemstyle{Auteur~:}
CNRM/GMEI/TRAMM

%%%%%REFERENCE%%%%
\itemstyle{R\'ef\'erences~:}
